%-------------------------
% Resume in LateX
% Author : Gabriel Gomez
% License : MIT
%------------------------

\documentclass[letterpaper,11pt]{article}

\usepackage{cmap}
\usepackage[utf8]{inputenc}
\usepackage[T1]{fontenc}
\usepackage{lmodern}

\usepackage{latexsym}
\usepackage[empty]{fullpage}
\usepackage{titlesec}
\usepackage{marvosym}
\usepackage[usenames,dvipsnames]{color}
\usepackage{verbatim}
\usepackage{enumitem}
\usepackage[hidelinks]{hyperref}
\usepackage{fancyhdr}
\usepackage[spanish]{babel}
\usepackage{tabularx}
\input{glyphtounicode}

\pagestyle{fancy}
\fancyhf{} % Clear all header and footer fields
\fancyfoot{}
\renewcommand{\headrulewidth}{0pt}
\renewcommand{\footrulewidth}{0pt}

% Adjust margins
\addtolength{\oddsidemargin}{-0.5in}
\addtolength{\evensidemargin}{-0.5in}
\addtolength{\textwidth}{1in}
\addtolength{\topmargin}{-.5in}
\addtolength{\textheight}{1.0in}

\urlstyle{same}

\raggedbottom
\raggedright
\setlength{\tabcolsep}{0in}

% Sections formatting
\titleformat{\section}{
  \vspace{-4pt}\scshape\raggedright\large
}{}{0em}{}[\color{black}\titlerule \vspace{-5pt}]

% Ensure that generate PDF is machine readable/ATS parsable
\pdfgentounicode=1

%-------------------------
% Custom commands
\newcommand{\resumeItem}[2]{
  \item\small{
    \textbf{#1}{: #2 \vspace{-2pt}}
  }
}

% Just in case someone needs a heading that does not need to be in a list
\newcommand{\resumeHeading}[4]{
    \begin{tabular*}{0.99\textwidth}[t]{l@{\extracolsep{\fill}}r}
      \textbf{#1} & #2 \\
      \textit{\small#3} & \textit{\small #4} \\
    \end{tabular*}\vspace{-5pt}
}

\newcommand{\resumeSubheading}[4]{
  \vspace{-1pt}\item
    \begin{tabular*}{0.97\textwidth}[t]{l@{\extracolsep{\fill}}r}
      \textbf{#1} & #2 \\
      \textit{\small#3} & \textit{\small #4} \\
    \end{tabular*}\vspace{-5pt}
}

\newcommand{\resumeSubSubheading}[2]{
    \begin{tabular*}{0.97\textwidth}{l@{\extracolsep{\fill}}r}
      \textit{\small#1} & \textit{\small #2} \\
    \end{tabular*}\vspace{-5pt}
}

\newcommand{\resumeSubItem}[2]{\resumeItem{#1}{#2}\vspace{-4pt}}

\renewcommand{\labelitemii}{$\circ$}

\newcommand{\resumeSubHeadingListStart}{\begin{itemize}[leftmargin=*]}
\newcommand{\resumeSubHeadingListEnd}{\end{itemize}}
\newcommand{\resumeItemListStart}{\begin{itemize}}
\newcommand{\resumeItemListEnd}{\end{itemize}\vspace{-5pt}}
\setlength{\footskip}{5pt}
%-------------------------------------------
%%%%%%  CV STARTS HERE  %%%%%%%%%%%%%%%%%%%%%%%%%%%%


\begin{document}

%----------HEADING-----------------
\begin{tabular*}{\textwidth}{l@{\extracolsep{\fill}}r}
  \textbf{\href{https://gabogomez.dev/}{\Large Gabriel Gomez}} & Correo electrónico: \href{mailto:ggomez.dev@outlook.com}{ggomez.dev@outlook.com}\\
  \href{https://gabogomez.dev/}{gabogomez.dev} & Movil: \href{tel:+573168325201}{+57-316-832-5201} \\
\end{tabular*}


%-----------EXPERIENCE-----------------
\section{Experiencia}
\resumeSubHeadingListStart

\resumeSubheading
{Autónomo}{Barranquilla, CO}
{Desarrollador de Software Freelance}{Jun 2025 -- Presente}
\resumeItemListStart
\resumeItem{Migración de Web Service}
{Migré un web service legacy de Java EE a Spring Boot aplicando Clean Architecture con separación en capas (dominio, aplicación, infraestructura) y principios SOLID. Implementé Kubernetes para la orquestación de contenedores y diseñé un entorno de desarrollo automatizado con Makefile que orquesta un stack Docker Compose con base de datos de pruebas, construcción de imágenes y variables de entorno por perfil.}
\resumeItem{Afinia PQR}
{Diseñé e implementé el backend Jakarta EE y la app Android para el sistema de gestión de PQR (peticiones, quejas y reclamos) de Afinia, como contratista freelance para Extreme Technologies. Modelé el dominio con una máquina de estados que controla las transiciones del caso a través de tres roles (administrador, contratista, supervisor) y diseñé las APIs RESTful que exponen el flujo de trabajo diferenciado por rol. Apliqué patrones GoF (State, Strategy) para encapsular la lógica de transición y mantener el backend extensible ante nuevos estados.}
\resumeItem{Sistema POS}
{Construí el backend con NestJS y SQLite para un sistema POS de escritorio offline-first, diseñando el esquema relacional con trazabilidad completa mediante tablas de auditoría para reducir discrepancias de inventario.}
\resumeItemListEnd

\resumeSubheading
{Extreme Technologies S.A.}{Barranquilla, CO}
{Desarrollador de Software}{Abr 2021 -- Jun 2024}
\resumeItemListStart
\resumeItem{Promigas}
{Diseñé e implementé el backend Jakarta EE de un sistema de gestión de cuadrillas para inspecciones de gasoductos, con rastreo GPS en tiempo real. Integré PostGIS para el procesamiento geoespacial de rutas y optimicé el web service para manejar sesiones concurrentes de alta carga, aplicando técnicas de concurrencia en Java para garantizar consistencia de datos bajo múltiples conexiones simultáneas.}
\resumeItem{Acuacar Daños}
{Mantuve y evoluccioné el backend Jakarta EE de un sistema de gestión de cuadrillas en la nube para Acuacar, implementando nuevas funcionalidades que abarcaron tanto el web service como la app móvil. Migré la arquitectura de la app para cumplir los requisitos de Android 11 y lideré el rediseño del flujo de trabajo.}
\resumeItem{Podas y Lavados}
{Desarrollé el backend Jakarta EE y la app Android multirol para el mantenimiento de redes eléctricas de Aire y Afinia, con soporte para dos flujos de trabajo paralelos. Implementé un módulo GIS que descarga y renderiza capas geoespaciales con el estado de postes y árboles desde el web service.}
\resumeItem{SETP Santa Marta}
{Implementé una capa middleware en Java sobre un validador de pasajes Android para transporte público, construyendo la lógica de validación de pagos (tarjetas físicas y QR) y diseñando una cola offline-first para almacenar transacciones durante pérdidas de conexión y sincronizarlas al restaurar la conectividad.}
\resumeItemListEnd

\resumeSubHeadingListEnd

%-----------EDUCATION-----------------
\section{Educación}
\resumeSubHeadingListStart
\resumeSubheading
{Universidad Autónoma del Caribe}{Barranquilla, CO}
{Pregrado en Ingeniería de Sistemas}{Feb 2018 -- May 2023}
\resumeSubHeadingListEnd


%-----------PROJECTS-----------------
\section{Proyectos}
\resumeSubHeadingListStart
\resumeSubItem{MatricuApp}
{Solución a prueba técnica — webapp de matrícula universitaria con autenticación JWT usando Spring Security, backend en Spring Boot con JPA y PostgreSQL (Neon), y frontend en React. Implementa control de créditos por semestre y gestión completa de materias.}
\resumeSubItem{Resume Builder on LaTeX}
{Generador de currículos de código abierto en LaTeX con GitHub Actions para compilación y despliegue automatizado de PDFs, implementando CI/CD con imágenes Docker multi-arquitectura publicadas en GHCR.}
\resumeSubItem{Portfolio Personal}
{Portfolio construido con Astro y Tailwind CSS, desplegado en Cloudflare con Edge Functions para redirección automática basada en idioma del navegador.}
\resumeSubHeadingListEnd

%-----------Cursos y Certificaciones-----------------
\section{Cursos y Certificaciones}
\resumeSubHeadingListStart
\resumeSubItem{Udemy - Curso de React: de cero a experto}{Mar 2026}
\resumeSubItem{Platzi - Herramientas de AI para Developers}{Abr 2025}
\resumeSubItem{Udemy - Jetpack Compose desde cero, migra tus vistas de Android XML}{Abr 2024}
\resumeSubItem{Udemy - Curso Testing para Android con JUnit, Mockito, Espresso y TDD}{Feb 2024}
\resumeSubItem{Universidad Autónoma del Caribe - Nuevas Tendencias de Ingeniería}{Feb 2023}
\resumeSubItem{Platzi - Curso Avanzado de Flutter}{Abr 2022}
\resumeSubItem{Platzi - Java SE Orientado a Objetos}{Feb 2021}
\resumeSubItem{Platzi - Introducción a Java SE}{Oct 2020}
\resumeSubHeadingListEnd

%--------SKILLS------------
\section{Skills}
\resumeSubHeadingListStart
\item{
      \textbf{Lenguajes}{: Java, Kotlin, TypeScript, SQL, JavaScript, HTML, CSS, Dart, Python, LaTeX}
}
\item{
      \textbf{Backend}{: Spring Boot, Jakarta EE, NestJS, RESTful APIs, JWT, Spring Security, Maven, Gradle}
}
\item{
      \textbf{Arquitectura}{: Clean Architecture, SOLID, Design Patterns (GoF), Hexagonal Architecture, Microservices, DDD, TDD, CI/CD, Agile, Scrum}
}
\item{
      \textbf{Datos}{: PostgreSQL, MySQL, SQLite, PostGIS, Firebase}
}
\item{
      \textbf{Infraestructura}{: AWS, Docker, Docker Compose, Kubernetes, GitHub Actions, Make}
}
\item{
      \textbf{Testing}{: JUnit, Mockito, Espresso, Test-Driven Development}
}
\resumeSubHeadingListEnd


%-------------------------------------------
\end{document}
